\section{Teori Informasi dan Entropi Shannon}

Keamanan sebuah cipher tidak hanya ditentukan oleh kerumitan algoritmanya, tetapi juga oleh seberapa banyak informasi yang "bocor" melalui cipherteks. Untuk mengukur hal ini secara objektif, kriptografi menggunakan \textbf{Teori Informasi}, yang dikembangkan oleh Claude Shannon pada tahun 1948.

\subsection{Konsep Entropi}

Dalam teori informasi, \textbf{entropi} adalah ukuran ketidakpastian atau jumlah rata-rata informasi yang terkandung dalam sebuah pesan. Secara intuitif, semakin acak sebuah data, semakin tinggi entropinya.

Jika kita memiliki variabel acak $X$ yang mewakili sebuah pesan dengan simbol-simbol $x_1, x_2, \dots, x_n$, dan setiap simbol memiliki peluang kemunculan $p(x_i)$, maka entropi $H(X)$ dihitung dengan rumus:
\[ H(X) = - \sum_{i=1}^{n} p(x_i) \log_2 p(x_i) \]
Satuan dari entropi adalah \textbf{bit per simbol}.

\textbf{Contoh Perhitungan}:
Misalkan kita memiliki pesan pendek dengan 8 karakter: \texttt{AABBCBDB}.
Simbol unik adalah $\{A, B, C, D\}$. Peluang kemunculannya:
$p(A) = 2/8, p(B) = 4/8, p(C) = 1/8, p(D) = 1/8$.

Maka entropinya adalah:
\begin{align*}
    H(X) &= - \left( \frac{2}{8} \log_2 \frac{2}{8} + \frac{4}{8} \log_2 \frac{4}{8} + \frac{1}{8} \log_2 \frac{1}{8} + \frac{1}{8} \log_2 \frac{1}{8} \right) \\
    &= - \left( 0.25 \cdot (-2) + 0.5 \cdot (-1) + 0.125 \cdot (-3) + 0.125 \cdot (-3) \right) \\
    &= - (-0.5 - 0.5 - 0.375 - 0.375) = 1.75 \text{ bit/simbol}
\end{align*}

\subsection{Rentang Entropi dan Signifikansinya}

Nilai entropi berada dalam rentang $0 \le H(X) \le \log_2(n)$, di mana $n$ adalah jumlah simbol unik.
\begin{itemize}
    \item \textbf{Entropi Minimum ($0$)}: Terjadi jika semua karakter dalam pesan identik (misal: \texttt{AAAAAAAA}). Tidak ada ketidakpastian sama sekali.
    \item \textbf{Entropi Maksimum ($\log_2 n$)}: Terjadi jika semua simbol muncul dengan peluang yang sama (distribusi seragam). Ini adalah kondisi yang paling acak.
\end{itemize}

Untuk teks bahasa Inggris yang terdiri dari 26 huruf alfabet, entropi maksimumnya adalah $\log_2 26 \approx 4.7$ bit/karakter. Namun, karena bahasa manusia memiliki pola (misal: huruf 'E' lebih sering muncul daripada 'Z'), entropi aktual bahasa Inggris jauh lebih rendah, sekitar $1.5$ bit/karakter.

\subsection{Hubungan Entropi dengan Keamanan Cipher}

Tujuan utama dari sebuah cipher yang baik adalah untuk memproduksi cipherteks yang memiliki \textbf{entropi maksimum}. Mengapa?
\begin{enumerate}
    \item \textbf{Menghilangkan Pola}: Jika cipherteks memiliki entropi rendah, maka pola-pola dari plainteks masih terbawa. Hal ini memudahkan kriptanalis melakukan analisis frekuensi.
    \item \textbf{Menyamar sebagai Noise}: Cipherteks dengan entropi tinggi akan terlihat seperti \textit{random noise} bagi pengamat luar. Penyerang tidak akan tahu apakah data tersebut adalah pesan rahasia atau sekadar sampah digital.
\end{enumerate}

Oleh karena itu, dalam desain cipher modern, digunakan operasi seperti S-Box (substitusi) dan P-Box (permutasi) untuk menyebarkan informasi sedemikian rupa sehingga distribusi karakter pada cipherteks mendekati distribusi seragam (entropi maksimal).
